% year: תשע"ז
% type: מבחן
% semester: א
% moed: א
% number: 6ב
% points: 10
% categories: improper-integrals, antiderivatives
% source: exams/מבחנים/תשעז/מועד א תשעז.pdf

\begin{question}
חשבו את האינטגרל הבא:
$\displaystyle\int_1^\infty x^2e^{-2x}\,dx$
\end{question}

\begin{hint}
בצעו אינטגרציה בחלקים פעמיים (גוזרים את $x^2$), ואז חשבו את הגבול כאשר הגבול העליון שואף לאינסוף.
\end{hint}

\begin{solution}
\textbf{פונקציה קדומה.} באינטגרציה בחלקים עם $u=x^2$, $dv=e^{-2x}dx$, $v=-\frac12e^{-2x}$:
\[
\int x^2e^{-2x}dx=-\frac{x^2}{2}e^{-2x}+\int xe^{-2x}dx.
\]
שוב בחלקים עם $u=x$:
\[
\int xe^{-2x}dx=-\frac x2e^{-2x}+\frac12\int e^{-2x}dx=-\frac x2e^{-2x}-\frac14e^{-2x}.
\]
לכן
\[
F(x)=\int x^2e^{-2x}dx=-e^{-2x}\left(\frac{x^2}{2}+\frac x2+\frac14\right)+C.
\]
\textbf{האינטגרל הלא אמיתי.} לפי ההגדרה,
\[
\int_1^\infty x^2e^{-2x}dx=\lim_{R\to\infty}\big(F(R)-F(1)\big).
\]
מתקיים $\lim_{R\to\infty}e^{-2R}\left(\frac{R^2}{2}+\frac R2+\frac14\right)=0$ (אקספוננט שולט בפולינום; למשל לפי כלל לופיטל פעמיים על $\frac{R^2}{e^{2R}}$), ו-$F(1)=-e^{-2}\left(\frac12+\frac12+\frac14\right)=-\frac54e^{-2}$. לכן
\[
\int_1^\infty x^2e^{-2x}dx=0+\frac54e^{-2}=\frac{5}{4e^2}\approx0.169.
\]
האינטגרל מתכנס, וערכו $\frac{5}{4e^2}$.
\end{solution}
